\chapter{The ordered bar construction}
\label{ch:ordered-bar}

\section{The simplicial object and the interval}

Let $A$ be a transversely ordered chiral algebra with augmentation
$\varepsilon\colon A\to\unit$ and augmentation ideal $\bar A$. The
simplicial ordered bar has level
\[
  B^{\mathrm{ord}}_r(A) = \unit\otimes^! A^{\otimes^! r}\otimes^!\unit,
\]
with faces built from the multiplication or the augmentation and
degeneracies built from the unit.

\begin{proposition}[The ordered bar factorises]
\label{prop:bar-factorises}
\StatusProved{} $B^{\mathrm{ord}}_\bullet(A)$ is a simplicial object of
$\FactD(X)$. Its geometric realization inherits the chiral factorisation
structure, and the deconcatenation coproduct is a morphism of
factorisation objects.
\end{proposition}

\begin{proof}
Faces and degeneracies are structure maps to which
Theorem~\ref{thm:locality} applies. Realization is computed
componentwise. Deconcatenation cuts a transverse word while a
factorisation equivalence separates the chiral support letter by letter;
the two commute on every disjointness locus, which is the morphism
condition in $\FactD(X)$.
\end{proof}

\begin{theorem}[Interval identity]
\label{thm:interval}
\StatusProved{} The realization of the simplicial ordered bar is the
relative tensor product of the two augmentation modules,
\[
  \Barperp_X(A) \;\simeq\; \unit\otimes^{L}_A\unit ,
\]
and when the analytic realization is locally constant in the transverse
direction and satisfies excision it is the factorisation homology of an
interval with augmentation boundary labels,
\[
  \Barperp_X(A) \;\simeq\; \int_{([0,1],\,\partial[0,1])} A .
\]
\end{theorem}

\begin{proof}
The first equivalence is the standard identification of the two-sided bar
resolution with the derived relative tensor product, valid in any
monoidal category with the relevant colimits; here the colimits are
computed componentwise on $X^I$ and remain in $\FactD(X)$ by
Proposition~\ref{prop:bar-factorises}. The second is factorisation
excision for the stratified interval, whose two boundary points carry the
augmentation modules, under the stated local-constancy hypothesis.
\end{proof}

\section{The chiral coefficient is not contracted}

The bar differential of Theorem~\ref{thm:interval} multiplies adjacent
transverse letters. It does not extract a collision residue. The de Rham
differential of a configuration space, the residue differential along a
collision divisor, and the simplicial bar differential are three distinct
operators, and each bar term carries the full chiral operation as a
morphism of factorisation $D$-modules.

\begin{remark}[The superposition that this separates]
\label{rem:superposition}
Placing a bar complex on $\overline{\mathcal M}_{g,n}$ over a relative
factorisation stack fuses the transverse and modular directions in the
primitive object. Every construction downstream of such a placement
inherits the fusion: one Maurer--Cartan equation where there are two, one
genus where there are two, and a single scalar where there are several
invariants in different theories. The separation is
Theorem~\ref{thm:interval} together with the independence of $g_\perp$
and $g_X$.
\end{remark}

\section{Coalgebra structure and completion}

The normalized complex of $\Barperp_X(A)$ is the conilpotent tensor
coalgebra $T^c(s\bar A)$ on the suspended augmentation ideal, with
deconcatenation coproduct. By Proposition~\ref{prop:bar-factorises} the
coproduct is a morphism of factorisation objects, so
\[
  \Barperp_X(A) \in
  \Coalg^{\mathrm{coaug}}_{E_1}\bigl(\FactD(X),\otimes^!\bigr)
\]
is coaugmented and conilpotent.

Two completions enter later statements and are never imposed silently:
the coradical filtration of the coalgebra, and the weight filtration when
$A$ is positively graded. Both are recorded in the hypotheses of
Chapter~\ref{ch:bar-cobar} rather than assumed here.

\begin{proposition}[Low arity]
\label{prop:low-arity}
\StatusProved{} In arity one the bar differential vanishes and
$\Barperp_X(A)_1 = s\bar A$. In arity two it is the suspended
multiplication $s\bar A\otimes^! s\bar A \to s\bar A$, and the
codimension-one boundary of the one-dimensional associahedron gives the
associativity relation $(ab)c-a(bc)=0$ in arity three.
\end{proposition}

\begin{proof}
Immediate from the simplicial identities, with the signs fixed by the
suspension convention: for a reduced bar word
$[a_1|\cdots|a_r]=sa_1\otimes\cdots\otimes sa_r$,
\[
  b_m[a_1|\cdots|a_r] = \sum_{i=1}^{r-1}
  (-1)^{\sum_{j\le i}|a_j| + i - 1}
  [a_1|\cdots|a_ia_{i+1}|\cdots|a_r].
\]
\end{proof}
